\section*{学习地图}
\addcontentsline{toc}{section}{学习地图}

\textbf{第 1、2 章}

\begin{center}
\footnotesize
\renewcommand{\arraystretch}{1.0}
\begin{tabularx}{\textwidth}{@{}>{\raggedright\arraybackslash}p{2.7cm}>{\raggedright\arraybackslash}p{4.1cm}>{\raggedright\arraybackslash}p{3.0cm}X@{}}
  \toprule
  单元 & 教材范围 & 学习日期 & 状态 \\
  \midrule
  \hyperref[section:QM-C01-S01]{QM-C01-S01} & 第 1 章 \S 1.1；印刷页 pp.~1--7；PDF pp.~15--21 & 2026-08-17 & 讲解、练习与收口检查完成 \\
  \hyperref[section:QM-C01-S02]{QM-C01-S02} & 第 1 章 \S 1.2；印刷页 pp.~7--11；PDF pp.~21--25 & 2026-08-18 & 讲解、同书练习与收口检查完成 \\
  \hyperref[section:QM-C01-S03]{QM-C01-S03} & 第 1 章 \S 1.3；印刷页 pp.~11--17；PDF pp.~25--31 & 2026-08-19 & 讲解、教材练习与收口检查完成 \\
  \hyperref[section:QM-C01-S04]{QM-C01-S04} & 第 1 章 \S 1.4；印刷页 pp.~17--18；PDF pp.~31--32 & 2026-08-20 & 讲解、教材练习与收口检查完成 \\
  \hyperref[section:QM-C01-S05]{QM-C01-S05} & 第 1 章 \S 1.5；印刷页 pp.~18--20；PDF pp.~32--34 & 2026-08-21 & 讲解、教材例子与自编收口检查完成 \\
  \hyperref[section:QM-C01-S06]{QM-C01-S06} & 第 1 章 \S 1.6；印刷页 pp.~20--29；PDF pp.~34--43 & 2026-08-22 & 讲解、教材习题与收口检查完成 \\
  \hyperref[section:QM-C01-S07]{QM-C01-S07} & 第 1 章 \S 1.7；印刷页 pp.~29--30；PDF pp.~43--44 & 2026-08-23 & 讲解、教材习题与收口检查完成 \\
  \hyperref[section:QM-C01-S08]{QM-C01-S08} & 第 1 章 \S 1.8；印刷页 pp.~30--54；PDF pp.~44--68 & 2026-08-24 至 2026-08-25 & 讲解、代表题与收口检查完成 \\
  \hyperref[section:QM-C01-S09]{QM-C01-S09} & 第 1 章 \S 1.9；印刷页 pp.~54--57；PDF pp.~68--71 & 2026-08-26 & 讲解、教材习题与收口检查完成 \\
  \hyperref[section:QM-C01-S10]{QM-C01-S10} & 第 1 章 \S 1.10；印刷页 pp.~57--73；PDF pp.~71--87 & 2026-08-26 至 2026-08-27 & 讲解、教材例题、习题与收口检查完成 \\
  \hyperref[section:QM-C02-S01]{QM-C02-S01} & 第 2 章 \S 2.1；印刷页 pp.~75--83；PDF pp.~88--96 & 2026-08-27 & 讲解、教材例题、习题与收口检查完成 \\
  \hyperref[section:QM-C02-S02]{QM-C02-S02} & 第 2 章 \S 2.2；印刷页 pp.~83--84；PDF pp.~96--97 & 2026-08-28 & 讲解、自编检验题与收口检查完成 \\
  \hyperref[section:QM-C02-S03]{QM-C02-S03} & 第 2 章 \S 2.3；印刷页 pp.~85--86；PDF pp.~98--99 & 2026-08-29 & 讲解、教材习题与收口检查完成 \\
  \hyperref[section:QM-C02-S04]{QM-C02-S04} & 第 2 章 \S 2.4；印刷页 p.~86；PDF p.~99 & 2026-08-30 & 讲解、自编检验题与收口检查完成 \\
  \hyperref[section:QM-C02-S05]{QM-C02-S05} & 第 2 章 \S 2.5；印刷页 pp.~86--90；PDF pp.~99--103 & 2026-08-31 & 讲解、教材例题、习题与收口检查完成 \\
  \hyperref[section:QM-C02-S06]{QM-C02-S06} & 第 2 章 \S 2.6；印刷页 pp.~90--91；PDF pp.~103--104 & 2026-09-01 & 讲解、自编检验题与收口检查完成 \\
  \hyperref[section:QM-C02-S07]{QM-C02-S07} & 第 2 章 \S 2.7；印刷页 pp.~91--98；PDF pp.~104--111 & 2026-09-02 & 讲解、教材习题与收口检查完成 \\
  \hyperref[section:QM-C02-S08]{QM-C02-S08} & 第 2 章 \S 2.8；印刷页 pp.~98--106；PDF pp.~111--119 & 2026-09-03 & 讲解、教材习题与收口检查完成 \\
  \bottomrule
\end{tabularx}
\end{center}

\textbf{主要教材：}Shankar, \textit{Principles of Quantum Mechanics}。

\textbf{记录规则：}区分教材内容、对话补充及自编题；解答见独立习题部分。

\clearpage
\section*{学习地图（续：第 3 至 7 章）}

\begin{center}
\footnotesize
\renewcommand{\arraystretch}{0.9}
\begin{tabularx}{\textwidth}{@{}>{\raggedright\arraybackslash}p{2.7cm}>{\raggedright\arraybackslash}p{4.1cm}>{\raggedright\arraybackslash}p{3.0cm}X@{}}
  \toprule
  单元 & 教材范围 & 学习日期 & 状态 \\
  \midrule
  \hyperref[section:QM-C03-S01]{QM-C03-S01} & 第 3 章 \S 3.1；印刷页 pp.~107--108；PDF pp.~120--121 & 2026-09-04 & 讲解、自编检查与订正完成 \\
  \hyperref[section:QM-C03-S02]{QM-C03-S02} & 第 3 章 \S 3.2；印刷页 pp.~108--110；PDF pp.~121--123 & 2026-09-05 & 讲解、自编检查与订正完成 \\
  \hyperref[section:QM-C03-S03]{QM-C03-S03} & 第 3 章 \S 3.3；印刷页 pp.~110--112；PDF pp.~123--125 & 2026-09-06 & 讲解、自编检查与订正完成 \\
  \hyperref[section:QM-C03-S04]{QM-C03-S04} & 第 3 章 \S 3.4；印刷页 p.~112；PDF p.~125 & 2026-09-07 & 讲解、教材估算例与自编检查完成 \\
  \hyperref[section:QM-C03-S05]{QM-C03-S05} & 第 3 章 \S 3.5；印刷页 pp.~112--113；PDF pp.~125--126 & 2026-09-08 & 讲解、自编检查与订正完成 \\
  \hyperref[section:QM-C04-S01]{QM-C04-S01} & 第 4 章 \S 4.1；印刷页 pp.~115--116；PDF pp.~127--128 & 2026-09-08 & 讲解、教材练习 4.2.1 与订正完成 \\
  \hyperref[section:QM-C04-S02]{QM-C04-S02} & 第 4 章 \S 4.2；印刷页 pp.~116--143；PDF pp.~128--155 & 2026-09-10 & 讲解、教材例题与练习 4.2.2--4.2.3 完成 \\
  \hyperref[section:QM-C04-S03]{QM-C04-S03} & 第 4 章 \S 4.3；印刷页 pp.~143--150；PDF pp.~155--162 & 2026-09-11 & 讲解、教材正文推导与应用、自编综合检验完成 \\
  \hyperref[section:QM-C05-S01]{QM-C05-S01} & 第 5 章 \S 5.1；印刷页 pp.~151--157；PDF pp.~163--169 & 2026-09-12 & 讲解、教材练习 5.1.1--5.1.4 与订正完成 \\
  \hyperref[section:QM-C05-S02]{QM-C05-S02} & 第 5 章 \S 5.2；印刷页 pp.~157--164；PDF pp.~169--176 & 2026-09-13 & 讲解、教材练习 5.2.1--5.2.6 与订正完成 \\
  \hyperref[section:QM-C05-S03]{QM-C05-S03} & 第 5 章 \S 5.3；印刷页 pp.~165--167；PDF pp.~177--179 & 2026-09-14 & 讲解、教材练习 5.3.1--5.3.4 与订正完成 \\
  \hyperref[section:QM-C05-S04]{QM-C05-S04} & 第 5 章 \S 5.4；印刷页 pp.~167--175；PDF pp.~179--187 & 2026-09-15 & 讲解、教材练习 5.4.1--5.4.3 与订正完成 \\
  \hyperref[section:QM-C05-S05]{QM-C05-S05} & 第 5 章 \S 5.5；印刷页 pp.~175--176；PDF pp.~187--188 & 2026-09-16 & 讲解、自编短应用题与核对完成 \\
  \hyperref[section:QM-C05-S06]{QM-C05-S06} & 第 5 章 \S 5.6；印刷页 pp.~176--178；PDF pp.~188--190 & 2026-09-17 & 两条定理、教材正文计算检验与核对完成 \\
  \hyperref[section:QM-C06-S01]{QM-C06-S01} & 第 6 章全文；印刷页 pp.~179--184；PDF pp.~191--196 & 2026-09-17 至 2026-09-18 & 讲解与例题完成；笔记已确认 \\
  \hyperref[section:QM-C07-S01]{QM-C07-S01} & 第 7 章 \S 7.1；印刷页 pp.~185--188；PDF pp.~197--200 & 2026-09-19 & 讲解与概念核对完成；笔记待确认 \\
  \bottomrule
\end{tabularx}
\end{center}
